\family{Malvaceae}{Mallow family}
\label{fam:malvaceae}

Trees, shrubs or herbs, often with star-shaped hairs. Leaves
alternate, simple or digitately compound, with stipules. Flowers
usually 5-merous, often large and showy. Fruit a capsule, a group of
follicles, or indehiscent. The family is taken here in its broad
modern sense, which includes the former Bombacaceae
(\taxon{Adansonia}, \taxon{Ceiba}) and Sterculiaceae (\taxon{Cola}).
Three species in the Arami Uplands.

\section*{Key to the species}
\begin{florakey}
	\lead{1a}{Leaves simple; fruit of up to 5 spreading follicles with
		large seeds}{\taxon{Cola nitida}}
	\lead{1b}{Leaves digitately compound}{2}
	\lead{2a}{Trunk massive and swollen; flowers large, white and pendent;
		fruit with a woody shell and a dry, powdery pulp}{\taxon{Adansonia
			digitata}}
	\lead{2b}{Trunk tall, with large buttresses and often spiny when young;
		fruit a capsule filled with silky floss}{\taxon{Ceiba pentandra}}
\end{florakey}

\species{Adansonia digitata}{L.}
\label{sp:MAL01}

\field{Local and trade names} Yoruba: Ose; Hausa: Kuka. Trade:
Baobab.

\field{Status} Native.

\field{Description} Deciduous tree to 20~m with a massive, swollen
trunk often more than 3~m in diameter. Bark smooth, grey. Leaves
digitately compound with 5--7 leaflets. Flowers large, white, pendent
on long stalks, opening at night. Fruit ovoid to oblong, 15--25~cm
long, with a velvety, woody shell and a whitish, powdery pulp.

\field{Phenology} Flowering April to June; fruiting October to
February.

\field{Habitat} Guinea savanna woodland, often near old settlements.

\field{Uses} The young leaves are cooked as a vegetable, the fruit
pulp is made into a drink, and the bark fibre is twisted into rope
\citep{burkill1985,arbonnier2004}.

\field{Specimens examined} Idofin (A3), 760~m, fl., 1~May 1961, \specimen{Olaniyan}{325}{SVN};
Okuta Hill (A4), 820~m, fr., 2~Dec 1962, \specimen{Olaniyan}{460}{SVN};
Aramiwo Escarpment, summit (B1), 900~m, fr., 15~Dec 1962, \specimen{Olaniyan}{481}{SVN};
Aramiwo Escarpment, summit (B1), 900~m, fr., 3~Nov 1967, \specimen{Olaniyan}{997}{SVN};
Asiri Crossing (A2), 740~m, st., 18~Sep 2020, \specimen{Alade}{327}{AUH};
Idofin (A3), 760~m, st., 10~Mar 2021, \specimen{Alade}{492}{AUH};
Aramiwo Escarpment, summit (B1), 900~m, fl., 23~May 2023, \specimen{Alade}{771}{AUH}.

\field{Distribution} Grid squares A2, A3, A4, B1.

\begin{center}
\distmap{A/2, A/3, A/4, B/1}
\end{center}

\species{Ceiba pentandra}{(L.) Gaertn.}
\label{sp:MAL02}

\field{Synonyms} \sciname{Bombax pentandrum}{L.}

\field{Local and trade names} Yoruba: Araba; Hausa: Rimi. Trade:
Ceiba.

\field{Status} Native.

\field{Description} Deciduous tree to 50~m with large plank
buttresses; young trunks and branches often spiny. Branches
horizontal, in tiers. Leaves digitately compound with 5--9 leaflets.
Flowers whitish to pale pink, appearing while the tree is leafless.
Fruit an oblong capsule 10--25~cm long, filled with silky floss
(kapok) round the seeds.

\field{Phenology} Flowering December to February; fruiting February
to April.

\field{Habitat} Lowland semi-deciduous forest and forest--savanna
mosaic; often left standing in villages and farmland.

\field{Uses} The floss is used to stuff mattresses and pillows, and
the light timber for canoes and boxes. Large trees are protected as
sacred in many villages \citep{burkill1985,keay1989}.

\field{Specimens examined} Ireke Ferry (D4), 360~m, fr., 3~Mar 1935, \specimen{Harding}{1353}{OKF};
Ireke Ferry (D4), 360~m, fl., 20~Feb 1936, \specimen{Harding}{1415}{OKF};
Okeri (C4), 420~m, fr., 6~Feb 1937, \specimen{Harding}{1616}{OKF};
Aramiwo Foot (B2), 640~m, fr., 16~Mar 1963, \specimen{Olaniyan}{607}{SVN};
Aramiwo Foot (B2), 640~m, fl., 7~Jan 1965, \specimen{Olaniyan}{697}{SVN};
Oke-Asa farmland (B4), 690~m, fr., 6~Mar 2019, \specimen{Alade}{185}{AUH};
Oke-Asa farmland (B4), 690~m, fr., 20~Apr 2023, \specimen{Alade}{719}{AUH}.

\field{Distribution} Grid squares B2, B4, C4, D4.

\begin{center}
\distmap{B/2, B/4, C/4, D/4}
\end{center}

\species{Cola nitida}{(Vent.) Schott \& Endl.}
\label{sp:MAL03}

\field{Synonyms} \sciname{Sterculia nitida}{Vent.}

\field{Local and trade names} Yoruba: Obi; Hausa: Goro. Trade: Kola.

\field{Status} Cultivated. Native to the forests further west and
widely planted in the Uplands.

\field{Description} Evergreen tree to 20~m. Leaves simple, elliptic
to obovate, 12--25~cm long, the petiole swollen at both ends. Flowers
cream with a purplish-red centre. Fruit of up to 5 spreading follicles,
each with several large red or white seeds (kola nuts).

\field{Phenology} Flowering March to June; fruiting October to
January.

\field{Habitat} Farms and compounds in the lowland semi-deciduous
forest and forest--savanna mosaic.

\field{Uses} The nuts contain caffeine and are chewed as a stimulant.
They are central to hospitality and ceremony and are traded widely
across West Africa \citep{keay1989,orwa2009}.

\field{Specimens examined} Ireke Ferry (D4), 360~m, fr., 9~Jan 1934, \specimen{Harding}{1049}{OKF};
Ireke Ferry (D4), 360~m, fr., 11~Nov 1934, \specimen{Harding}{1236}{OKF};
Olowo Stream (D3), 440~m, fl., 18~Jun 1938, \specimen{Harding}{1970}{OKF};
Aramiwo Foot (B2), 640~m, fr., 2~Jan 1961, \specimen{Olaniyan}{247}{SVN};
Aramiwo Foot (B2), 640~m, st., 10~Aug 1962, \specimen{Olaniyan}{450}{SVN};
Aramiwo Foot (B2), 640~m, fl., 12~Jun 2020, \specimen{Alade}{282}{AUH};
Agbede (C2), 540~m, st., 2~Aug 2020, \specimen{Alade}{297}{AUH}.

\field{Distribution} Grid squares B2, C2, D3, D4.

\begin{center}
\distmap{B/2, C/2, D/3, D/4}
\end{center}
